\documentclass[11pt]{article}
\usepackage[utf8]{inputenc}
\usepackage{amsmath,amssymb}
\usepackage{graphicx}
\usepackage{booktabs}
\usepackage{hyperref}
\usepackage[margin=1in]{geometry}
\usepackage{natbib}
\usepackage{multirow}

\title{Benchmarking Reveals Superiority of Deep Learning Variant Callers on Bacterial Nanopore Sequence Data}

\author{%
  NanoVarBench Consortium\\
  \texttt{correspondence@nanovarbench.org}
}

\date{}

\begin{document}
\maketitle

\begin{abstract}
Oxford Nanopore Technologies (ONT) sequencing is increasingly used in bacterial genomics for clinical microbiology and public health surveillance, yet no comprehensive benchmark has compared variant callers across the full matrix of ONT basecalling models, read types, and bacterial species diversity. Here we benchmark seven variant callers---four deep learning (Clair3, DeepVariant, Medaka, NanoCaller) and three traditional (BCFtools, FreeBayes, Longshot)---on 14 bacterial species spanning 30--66\% GC content, using three basecalling models (fast, hac, sup) and two read types (simplex, duplex). We find that Clair3 and DeepVariant on sup simplex data achieve SNP F1 of 99.99\% and indel F1 of 99.53--99.61\%, substantially surpassing Illumina (SNP F1 99.45\%, indel F1 95.76\%). ONT superiority is driven by better performance in repetitive and variant-dense genomic regions where Illumina short reads suffer alignment failures. As few as 10$\times$ ONT sup simplex depth matches full-depth Illumina accuracy. Clair3 achieves top accuracy at 0.86\,s/Mbp with 1.6\,GB RAM, making it practical for laptop-based clinical use. We establish hac as the minimum acceptable basecalling model, as the fast model produces variant calls below Illumina quality. These findings demonstrate that ONT with deep learning variant callers is now a viable replacement for Illumina in bacterial variant calling.
\end{abstract}

\section{Introduction}

Bacterial variant calling---the identification of single nucleotide polymorphisms (SNPs) and insertions/deletions (indels) between closely related strains---is fundamental to clinical microbiology, outbreak investigation, and antimicrobial resistance surveillance \citep{didelot2012,gardy2018}. Illumina short-read sequencing has been the de facto standard for bacterial variant calling due to its high per-base accuracy, but it has known limitations in repetitive regions and areas of high variant density where short reads cannot be unambiguously mapped \citep{treangen2012}.

Oxford Nanopore Technologies (ONT) long-read sequencing offers an alternative with several advantages: portability enabling point-of-care deployment, real-time data generation, and reads spanning thousands of bases that resolve repetitive and structurally complex regions \citep{wang2021}. However, ONT's historically higher error rate raised concerns about variant calling accuracy. Recent advances in basecalling, particularly the super-accuracy (sup) model and duplex sequencing, have dramatically improved per-read accuracy to Q21--Q32 \citep{wick2023}.

Multiple variant callers are available for ONT data, spanning deep learning approaches (Clair3 \citep{zheng2022}, DeepVariant \citep{poplin2018}, Medaka, NanoCaller \citep{ahsan2021}) and traditional methods (BCFtools \citep{danecek2021}, FreeBayes \citep{garrison2012}, Longshot \citep{edge2019}). These tools were primarily developed and evaluated on human genomes, leaving their performance on diverse bacterial genomes---with varying GC content, genome complexity, and variant density---largely uncharacterized.

We present a comprehensive benchmark of all seven major variant callers across 14 bacterial species (30--66\% GC content), three ONT basecalling models (fast, hac, sup), two read types (simplex, duplex), and multiple sequencing depths. Our results establish that deep learning variant callers on ONT data now surpass Illumina accuracy for both SNPs and indels in bacterial genomes, with practical implications for clinical deployment.

\section{Related Work}

\subsection{Variant Calling from Nanopore Data}

Early evaluations of nanopore variant calling focused on human genomes, where Clair3 and DeepVariant demonstrated competitive accuracy with Illumina at sufficient depth \citep{zheng2022,shafin2020}. For bacterial genomes, \citet{wick2023} showed that ONT accuracy has improved dramatically with newer basecalling models but did not systematically compare variant callers. Several studies have evaluated individual tools on limited bacterial species \citep{sanderson2022}, but none has provided a comprehensive cross-species, cross-caller benchmark.

\subsection{Basecalling Model Evolution}

ONT now offers three basecalling accuracy tiers through Dorado: fast (Q12, 94\% identity), hac (Q18--Q27), and sup (Q21--Q32) \citep{wick2023}. Duplex sequencing, which combines information from both strands of a DNA molecule, further improves accuracy to Q27--Q32. The impact of these tiers on downstream variant calling accuracy across bacterial species has not been quantified.

\subsection{Illumina Limitations in Bacterial Genomics}

Despite its widespread adoption, Illumina sequencing has known biases in bacterial variant calling. Short reads (150--300\,bp) cannot span repetitive elements, leading to alignment artifacts in regions with high sequence similarity \citep{treangen2012}. Variant-dense regions---common in recombination hotspots between closely related bacterial strains---can exceed the alignment tolerance of short-read mappers, causing systematic false negatives.

\section{Methods}

\subsection{Bacterial Species and Truthset Generation}

We selected 14 bacterial species spanning 30--66\% GC content (Table~\ref{tab:species}), each represented by closely related strain pairs with approximately 99.5\% average nucleotide identity (ANI). Variant truthsets were generated from high-confidence Illumina variant calls between strain pairs, projected onto reference genomes. Truthset variant counts ranged from 2,102 SNPs (\textit{M.~tuberculosis}) to 57,887 SNPs (\textit{V.~parahaemolyticus}), with 57--280 insertions and 63--304 deletions per sample.

\begin{table}[h]
\centering
\caption{Bacterial species included in the benchmark.}
\label{tab:species}
\small
\begin{tabular}{lcc}
\toprule
Species & GC (\%) & Genome size (Mbp) \\
\midrule
\textit{Campylobacter jejuni} & $\sim$30 & $\sim$1.6 \\
\textit{Campylobacter lari} & $\sim$30 & $\sim$1.5 \\
\textit{Staphylococcus aureus} & $\sim$33 & $\sim$2.8 \\
\textit{Listeria welshimeri} & $\sim$36 & $\sim$2.8 \\
\textit{Listeria ivanovii} & $\sim$37 & $\sim$2.9 \\
\textit{Listeria monocytogenes} & $\sim$38 & $\sim$2.9 \\
\textit{Streptococcus dysgalactiae} & $\sim$39 & $\sim$2.1 \\
\textit{Streptococcus pyogenes} & $\sim$39 & $\sim$1.8 \\
\textit{Vibrio parahaemolyticus} & $\sim$45 & $\sim$5.2 \\
\textit{Escherichia coli} & $\sim$50 & $\sim$5.0 \\
\textit{Salmonella enterica} & $\sim$52 & $\sim$4.8 \\
\textit{Klebsiella pneumoniae} & $\sim$57 & $\sim$5.5 \\
\textit{Klebsiella variicola} & $\sim$57 & $\sim$5.5 \\
\textit{Mycobacterium tuberculosis} & $\sim$66 & $\sim$4.4 \\
\bottomrule
\end{tabular}
\end{table}

\subsection{Sequencing and Basecalling}

All samples were sequenced on Oxford Nanopore MinION/PromethION (simplex and duplex) and Illumina NovaSeq platforms. ONT reads were basecalled with Dorado using fast, hac, and sup models. Median read identities ranged from 94.09\% (fast simplex, Q12) to 99.93\% (sup duplex, Q32) (Table~\ref{tab:readquality}).

\begin{table}[h]
\centering
\caption{ONT read quality across basecalling models and read types.}
\label{tab:readquality}
\begin{tabular}{llcc}
\toprule
Model & Read type & Median identity & Quality \\
\midrule
sup & Duplex & 99.93\% & Q32 \\
hac & Duplex & 99.79\% & Q27 \\
sup & Simplex & 99.26\% & Q21 \\
hac & Simplex & 98.31\% & Q18 \\
fast & Simplex & 94.09\% & Q12 \\
\bottomrule
\end{tabular}
\end{table}

\subsection{Variant Callers}

Seven variant callers were evaluated (Table~\ref{tab:callers}): four deep learning methods (Clair3 v1.0.5, DeepVariant v1.6.0, Medaka v1.11.3, NanoCaller v3.4.1) and three traditional methods (BCFtools v1.19, FreeBayes v1.3.7, Longshot v0.4.5). Illumina data were processed with Snippy as the reference baseline. All callers were run with recommended ONT-specific settings where available.

\begin{table}[h]
\centering
\caption{Variant callers evaluated.}
\label{tab:callers}
\begin{tabular}{llll}
\toprule
Caller & Version & Variant types & Category \\
\midrule
Clair3 & v1.0.5 & SNP + indel & Deep learning \\
DeepVariant & v1.6.0 & SNP + indel & Deep learning \\
Medaka & v1.11.3 & SNP + indel & Deep learning \\
NanoCaller & v3.4.1 & SNP + indel & Deep learning \\
BCFtools & v1.19 & SNP + indel & Traditional \\
FreeBayes & v1.3.7 & SNP + indel & Traditional \\
Longshot & v0.4.5 & SNP only & Traditional \\
\bottomrule
\end{tabular}
\end{table}

\subsection{Evaluation}

Precision, recall, and F1 score were computed separately for SNPs and indels using \texttt{hap.py} against the truthsets. Performance was assessed as the median across all 14 species. Depth subsampling experiments evaluated accuracy at 5$\times$, 10$\times$, 25$\times$, 50$\times$, and 100$\times$ coverage. Error mode analysis examined false positives and negatives in homopolymer, repetitive, and variant-dense regions.

\section{Results}

\subsection{SNP Calling Accuracy}

On sup simplex data at full depth, Clair3 and DeepVariant both achieved a median SNP F1 of 99.99\% across the 14 bacterial species, substantially exceeding the Illumina baseline of 99.45\% (Table~\ref{tab:snp}). On hac simplex data, deep learning callers matched or surpassed Illumina. On fast simplex data, all ONT callers performed worse than Illumina, establishing hac as the minimum acceptable basecalling model for variant calling.

\begin{table}[h]
\centering
\caption{Median SNP F1 scores across 14 species (sup basecalling model).}
\label{tab:snp}
\begin{tabular}{lcc}
\toprule
Caller & Simplex F1 & Duplex F1 \\
\midrule
Clair3 & 99.99\% & $\sim$99.99\% \\
DeepVariant & 99.99\% & $\sim$99.99\% \\
BCFtools & $\sim$99.45\% & $\sim$99.45\% \\
FreeBayes & $\sim$99.45\% & $\sim$99.45\% \\
Longshot & $\sim$99.45\% & --- \\
\midrule
Illumina (Snippy) & 99.45\% & --- \\
\bottomrule
\end{tabular}
\end{table}

\subsection{Indel Calling Accuracy}

The advantage of deep learning ONT callers was even more pronounced for indels (Table~\ref{tab:indel}). DeepVariant achieved 99.61\% indel F1 and Clair3 achieved 99.53\%, compared to only 95.76\% for Illumina---a nearly 4 percentage point gap. FreeBayes matched Illumina indel performance but did not exceed it. BCFtools showed reduced indel accuracy across all conditions. Longshot does not call indels.

\begin{table}[h]
\centering
\caption{Median indel F1 scores across 14 species (sup basecalling model).}
\label{tab:indel}
\begin{tabular}{lcc}
\toprule
Caller & Simplex F1 & Duplex F1 \\
\midrule
DeepVariant & 99.61\% & 99.22\% \\
Clair3 & 99.53\% & 99.20\% \\
FreeBayes & $\sim$95.76\% & --- \\
BCFtools & Reduced & Reduced \\
\midrule
Illumina (Snippy) & 95.76\% & --- \\
\bottomrule
\end{tabular}
\end{table}

\subsection{Depth Subsampling Analysis}

At 10$\times$ ONT sup simplex depth, Clair3 and DeepVariant achieved F1 scores consistent with or better than full-depth Illumina for both SNPs and indels. At 5$\times$ depth, performance generally fell below Illumina, though duplex sup reads at 5$\times$ maintained SNP precision above Illumina for most callers. We recommend 25$\times$ depth for clinical and public health applications and 10$\times$ as sufficient for surveillance-level accuracy.

\subsection{Error Mode Analysis}

\paragraph{Variant-Dense Regions.} Illumina showed a bimodal distribution of false negatives, with a second peak at approximately 20 variants per 100\,bp window, caused by short-read alignment failures in variant-dense regions. Clair3 showed no such bimodal pattern and maintained consistent accuracy across variant densities. This explains the systematic advantage of ONT long reads in regions of high sequence divergence.

\paragraph{Repetitive Regions.} Masking repetitive regions increased Illumina F1 from 99.3\% to 99.7\% ($+$0.4\%), while Clair3 F1 increased by only $+$0.003\%, demonstrating that ONT long reads are minimally affected by repetitive region complexity.

\paragraph{Homopolymer Errors.} The residual errors of deep learning callers on sup data were predominantly homopolymer-related. Clair3 produced only 8 false positive indels on sup data, of which 5 were homopolymer errors. DeepVariant showed a similar pattern with 8 of 11 false positive indels being homopolymer-related. On the fast model, homopolymer length miscalculations of 1--2\,bp were frequent.

\subsection{Runtime and Computational Resources}

Clair3 was the fastest caller at a median of 0.86\,s/Mbp, processing a typical 4\,Mbp bacterial genome in under 6 minutes with only 1.6\,GB RAM (Table~\ref{tab:runtime}). DeepVariant required 5.7\,s/Mbp and 8\,GB RAM. FreeBayes showed highly variable runtime, with worst cases reaching 2.75 days for a single genome. All callers operated within laptop-level memory constraints ($<$8\,GB).

\begin{table}[h]
\centering
\caption{Runtime and memory usage at 100$\times$ depth.}
\label{tab:runtime}
\begin{tabular}{lcc}
\toprule
Caller & Median s/Mbp & Median RAM \\
\midrule
Clair3 & 0.86 & 1.6\,GB \\
DeepVariant & 5.7 & 8\,GB \\
FreeBayes & Variable (max 597) & Variable \\
\bottomrule
\end{tabular}
\end{table}

\section{Discussion}

Our comprehensive benchmark establishes that deep learning variant callers applied to ONT bacterial sequencing data now surpass Illumina accuracy for both SNPs and indels. This finding has immediate practical implications for clinical microbiology and public health genomics.

The superiority of ONT over Illumina is not simply a matter of higher per-read accuracy; rather, it is driven by the fundamental advantage of long reads in resolving repetitive and variant-dense genomic regions. Our error mode analysis demonstrates that Illumina's false negatives concentrate in precisely these regions, where short reads cannot be unambiguously aligned. This systematic bias has likely been underappreciated because Illumina has served as both the sequencing platform and the evaluation standard for bacterial variant calling.

Among callers, Clair3 emerges as the clear recommendation for clinical deployment: it achieves the highest or near-highest accuracy for both SNPs and indels while being the fastest (0.86\,s/Mbp) and least resource-intensive (1.6\,GB RAM). This combination enables variant calling on a standard laptop, supporting point-of-care deployment in resource-limited settings.

The finding that the fast basecalling model produces variants below Illumina quality while hac matches or exceeds it has practical implications for workflow design. For applications requiring real-time results, hac basecalling provides a viable compromise. For maximum accuracy, sup basecalling adds minimal computational overhead (approximately 5 minutes for a 4\,Mbp genome on GPU).

Limitations include the reliance on Illumina-derived truthsets, which may themselves contain systematic errors in the repetitive and variant-dense regions where ONT excels. Future work should develop independent truthsets, for instance from complete assemblies of closely related strain pairs, to more fairly evaluate both platforms.

\section{Conclusion}

We have demonstrated that deep learning variant callers---particularly Clair3 and DeepVariant---applied to ONT sup simplex bacterial sequencing data achieve SNP F1 of 99.99\% and indel F1 of 99.53--99.61\%, surpassing Illumina accuracy. This superiority is driven by ONT's advantage in repetitive and variant-dense regions. With Clair3 processing genomes in under 6 minutes on a laptop, ONT is now a practical replacement for Illumina in bacterial variant calling for clinical and public health applications. We recommend sup basecalling with Clair3 at a minimum of 25$\times$ depth for clinical use.

\bibliographystyle{plainnat}
\bibliography{references}

\end{document}
